
\section{Formal Analysis}
\label{sec:formal-analysis}

This analysis characterizes how explicit project state affects
project-continuation decisions.  Code quality remains an empirical question
outside the scope of the formal results.  The main results connect
action-relevant observation aliasing to unavoidable continuation error or
evidence-recovery cost, then identify when an explicit state projection is
sufficient to remove hidden coordination-state aliasing.  Transactional
soundness ensures that the projection comes from one coherent committed state.

\subsection{Observation Aliasing and Continuation Error}

Our claims concern \emph{project-level continuation decisions} governed by the
admissible transition relation $\Gamma$, rather than arbitrary local code edits.
Let $\Omega$ be the set of project worlds.  Each $\omega\in\Omega$ induces a
repository $X(\omega)$, a historical memory collection $H(\omega)$, and an
accepted coordination state $z(\omega)\in\Z$.  For a fixed continuation
query $q$, define the observation available under repository budget $\beta$ and
retrieval depth $k$ as
\begin{equation}
\phi_{\beta,k}(q,\omega)
=
\Bigl(
q,\,
O_\beta(q,X(\omega)),\,
R_k(q,H(\omega))
\Bigr).
\label{eq:observation-map}
\end{equation}

Let $\mathcal U_c$ be the set of project-level continuation actions, such as
entering, pausing, aborting, or completing a branch.  Given
$\Gamma\subseteq\mathcal Z\times\mathcal U_c\times\mathcal Z$, define the
admissible continuation set
\begin{equation}
\operatorname{Adm}_\Gamma(z)
=
\left\{
u\in\mathcal U_c:
\exists z'\in\mathcal Z,\ (z,u,z')\in\Gamma
\right\}.
\label{eq:admissible-set}
\end{equation}

\begin{definition}[Decisive observation alias]
Two project worlds $\omega_1,\omega_2$ are decisively aliased under
$\phi_{\beta,k}$ for query $q$ if
\begin{equation}
\phi_{\beta,k}(q,\omega_1)
=
\phi_{\beta,k}(q,\omega_2)
\label{eq:same-observation}
\end{equation}
while
\begin{equation}
\operatorname{Adm}_\Gamma(z(\omega_1))
\cap
\operatorname{Adm}_\Gamma(z(\omega_2))
=
\varnothing.
\label{eq:disjoint-actions}
\end{equation}
Thus, the available evidence is identical even though no project-level
continuation action is valid in both worlds.
\end{definition}

\begin{theorem}[Continuation error under decisive aliasing]
\label{thm:aliasing-lower-bound}
Let $\omega_1,\omega_2$ be decisively aliased, and let
$\pi(\cdot\mid o)$ be any possibly randomized policy that chooses a continuation
action using only the observation
$o=\phi_{\beta,k}(q,\omega_i)$.  Under prior probabilities
$p$ and $1-p$ for the two worlds, respectively,
\begin{equation}
\Pr[\textnormal{admissible continuation}]
\le
\max\{p,1-p\}.
\label{eq:aliasing-success-bound}
\end{equation}
Equivalently,
\begin{equation}
\Pr[\textnormal{continuation error}]
\ge
\min\{p,1-p\}.
\label{eq:aliasing-error-bound}
\end{equation}
In particular, under an equal prior, every observation-only policy has
continuation error at least $1/2$ on the pair.
\end{theorem}

\begin{proof}
Write
$A_i=\operatorname{Adm}_\Gamma(z(\omega_i))$.
Because the observations are identical, the policy uses the same distribution
$\pi(\cdot\mid o)$ in both worlds.  Define
$\pi(A\mid o)=\sum_{u\in A}\pi(u\mid o)$.  Its success probability is
\begin{equation}
p\,\pi(A_1\mid o)
+
(1-p)\,\pi(A_2\mid o).
\end{equation}
Since $A_1\cap A_2=\varnothing$,
$\pi(A_1\mid o)+\pi(A_2\mid o)\le 1$.  Therefore
\begin{align}
p\,\pi(A_1\mid o)
+
(1-p)\,\pi(A_2\mid o)
&\le
\max\{p,1-p\}
\bigl(
\pi(A_1\mid o)+\pi(A_2\mid o)
\bigr)\\
&\le
\max\{p,1-p\}.
\end{align}
Taking one minus both sides yields
Equation~\eqref{eq:aliasing-error-bound}.
\end{proof}

\begin{corollary}[Repository-only non-identifiability]
\label{cor:repository-nonidentifiability}
Suppose two distinct development histories satisfy
\begin{equation}
\Psi(h,\alpha)=\Psi(h',\alpha')=X
\end{equation}
but induce disjoint admissible continuation sets.  Then any policy whose input is
only a function of the repository snapshot $X$ is subject to
Theorem~\ref{thm:aliasing-lower-bound}.
\end{corollary}

\begin{proof}
Every code-only observation has the same value in the two worlds because the
repository snapshot is identical.  The worlds are therefore decisively aliased,
and Theorem~\ref{thm:aliasing-lower-bound} applies.
\end{proof}

\paragraph{A TreeWork witness.}
The condition above can arise from branch-bound execution rather than from an
artificial complete-versus-incomplete pair.  Let $b_1\neq b_2$ be two
nonterminal branches, and let $q$ ask an agent to continue the branch assigned
to its current bound workspace.  Consider two worlds with the same observed
source tree and the same retrieved historical records, while the hidden
worktree binding resolves the invocation to $b_1$ in $\omega_1$ and to $b_2$ in
$\omega_2$.  Instantiate the query-specific continuation action set as
\begin{equation}
\mathcal U_c(q)=\{\operatorname{Enter}(b_1),\operatorname{Enter}(b_2)\}.
\end{equation}
TreeWork's binding guard rejects cross-branch entry, so the admissible sets are
$\{\operatorname{Enter}(b_1)\}$ and $\{\operatorname{Enter}(b_2)\}$,
respectively.  They are disjoint even though $\phi_{\beta,k}$ is identical when
the binding is absent from the repository and retrieval observations.  This
witness establishes realizability of decisive aliasing in the implemented
coordination model; it does not assert a population frequency for such pairs.

\subsection{Recovery Cost}

An additional evidence-acquisition procedure may inspect more files, search Git
history, retrieve more memories, run tools, or request human clarification.  We
model $\sigma$ as a deterministic adaptive policy that maps its current
transcript to either the next evidence request or $\operatorname{Stop}$.  Let
$T_\sigma(\omega)$ be its terminal transcript in world $\omega$, and let
$c_\sigma(\omega)\ge 0$ be the accumulated token, tool, time, or
human-interaction cost in that world.  For a pair of worlds, define the minimum
worst-case separation cost
\begin{equation}
C_{\mathrm{sep}}(\omega_1,\omega_2)
=
\inf_{\sigma:
T_\sigma(\omega_1)\neq T_\sigma(\omega_2)}
\max_{i\in\{1,2\}}c_\sigma(\omega_i),
\label{eq:separation-cost}
\end{equation}
with $C_{\mathrm{sep}}=+\infty$ when no available procedure distinguishes the
pair.

\begin{theorem}[Error--recovery-cost dichotomy]
\label{thm:error-or-cost}
For a decisively aliased pair $\omega_1,\omega_2$, any deterministic adaptive
continuation procedure that guarantees an admissible action in both worlds must
incur worst-case additional evidence-acquisition cost at least
$C_{\mathrm{sep}}(\omega_1,\omega_2)$.  If
$C_{\mathrm{sep}}(\omega_1,\omega_2)=+\infty$, no zero-error continuation
procedure exists.
\end{theorem}

\begin{proof}
Suppose a procedure guarantees an admissible action in both worlds.  Its
terminal transcripts must differ: identical transcripts cause the deterministic
decision rule to choose the same action, which cannot belong to two disjoint
admissible sets.  The procedure is therefore among those in the infimum in
Equation~\eqref{eq:separation-cost}, and its worst-case cost is at least
$C_{\mathrm{sep}}$.  If $C_{\mathrm{sep}}=+\infty$, no available procedure
produces distinct terminal transcripts, so zero-error continuation is
impossible.
\end{proof}

\subsection{Continuation Sufficiency of TreeWork}

For each continuation action $u\in\mathcal U_c$, let
\begin{equation}
g_u(z)
=
\mathbf 1
\left[
u\in\operatorname{Adm}_\Gamma(z)
\right]
\end{equation}
be its admissibility guard.

\begin{definition}[Guard-complete projection]
A projection $\Pi$ of maintained project memory is guard-complete for
$\mathcal U_c$ if, for every $u\in\mathcal U_c$, there exists a deterministic
function $\widetilde g_u$ such that
\begin{equation}
g_u(z(M))
=
\widetilde g_u(\Pi(M))
\qquad
\textnormal{for every committed project memory }M.
\label{eq:guard-factorization}
\end{equation}
\end{definition}

\begin{lemma}[Guard agreement from a guard-complete projection]
\label{thm:continuation-sufficiency}
If $\Pi$ is guard-complete for $\mathcal U_c$, then
\begin{equation}
\Pi(M)=\Pi(M')
\quad\Longrightarrow\quad
\operatorname{Adm}_\Gamma(z(M))
=
\operatorname{Adm}_\Gamma(z(M')).
\label{eq:sufficient-projection}
\end{equation}
Equivalently, two states with different admissible continuation sets cannot
share the same projection.
\end{lemma}

\begin{proof}
Assume $\Pi(M)=\Pi(M')$.  For every $u\in\mathcal U_c$,
Equation~\eqref{eq:guard-factorization} gives
\begin{equation}
g_u(z(M))
=
\widetilde g_u(\Pi(M))
=
\widetilde g_u(\Pi(M'))
=
g_u(z(M')).
\end{equation}
Thus exactly the same actions satisfy their guards in both states, which proves
Equation~\eqref{eq:sufficient-projection}.
\end{proof}

This lemma records a representational consequence of guard completeness; it is
not by itself a TreeWork-specific guarantee.  The substantive obligation is to
show that a concrete projection contains every field consulted by the runtime
guards.

\begin{proposition}[TreeWork hard-guard projection is continuation-sufficient]
\label{prop:treework-guard-complete}
For the branch-lifecycle commands used during Work Tree, fix an invocation
context $j$.  Let
\begin{equation}
\begin{aligned}
\mathcal U_b=\{&\operatorname{Enter}(b),\operatorname{Pause}(b),\\
                &\operatorname{Abort}(b),\operatorname{Complete}(b)\}.
\end{aligned}
\end{equation}
Let $\operatorname{Targetable}_t(j,b)$ indicate that a branch-scoped command
issued from $j$ resolves to $b$, and let
$\operatorname{EnterSafe}_t(j,b)$ summarize the runtime's worktree-binding and
isolation checks for entering $b$.  Define the hard-guard projection
\begin{equation}
\begin{aligned}
\chi_t(b,j)
=
\Bigl(&\ell_t(b),v_t(b),
  \operatorname{Targetable}_t(j,b),\\
  &\operatorname{EnterSafe}_t(j,b),
  \operatorname{AcceptanceDone}_t(b),\\
  &\operatorname{StateConsistent}_t,
  \operatorname{WorkspaceSafe}_t(b)\Bigr).
\end{aligned}
\label{eq:coordination-projection}
\end{equation}
Assume the environment predicates report exactly the checks performed by the
runtime at transaction time.  Then $\chi_t(b,j)$ is guard-complete for the
hard lifecycle and completion guards over $\mathcal U_b$.
\end{proposition}

\begin{proof}
The guard for $\operatorname{Enter}(b)$ depends on a nonterminal lifecycle and
the invocation and isolation checks summarized by
$\operatorname{Targetable}_t(j,b)$ and $\operatorname{EnterSafe}_t(j,b)$.
The guards for $\operatorname{Pause}(b)$ and $\operatorname{Abort}(b)$ depend on
target resolution and lifecycle.  The guard for
$\operatorname{Complete}(b)$ additionally depends on
$\operatorname{AcceptanceDone}_t(b)$, $v_t(b)$,
$\operatorname{StateConsistent}_t$, and
$\operatorname{WorkspaceSafe}_t(b)$.  Every hard guard therefore factors
through $\chi_t(b,j)$, so
Lemma~\ref{thm:continuation-sufficiency} applies.
\end{proof}

Table~\ref{tab:runtime-correspondence} maps the predicates in $\chi_t$ to the
runtime mechanisms and regression evidence used by the implementation.  The
proposition remains conditional on those mechanisms reporting the checks
performed at transaction time.

The runtime serializes an Agent-visible Recall interface containing the
machine-derived allowed-action set, blocked actions with reason codes, and a
publication marker $m_t=(r_t,e_t,h_t)$ containing the Tree revision, event
sequence, and accepted Tree hash.  Let $D_{\mathrm{allow}}$ decode the serialized
allowed-action list.

\begin{proposition}[Recall interface correspondence]
\label{prop:recall-interface-correspondence}
Assume Recall eligibility and lifecycle commands evaluate the hard guards in
Proposition~\ref{prop:treework-guard-complete} over the same committed state and
invocation context.  Then, at read time,
\begin{equation}
D_{\mathrm{allow}}
\bigl(\operatorname{RecallAPI}_t(b,j)\bigr)
=
\operatorname{Adm}_\Gamma(z(M_t))\cap\mathcal U_b.
\label{eq:recall-interface-correspondence}
\end{equation}
\end{proposition}

\begin{proof}
Recall partitions every action in $\mathcal U_b$ into allowed or blocked using
the same lifecycle, targetability, acceptance, verification, consistency, and
workspace predicates represented by $\chi_t(b,j)$.  The decoder selects exactly
the allowed partition.  Guard completeness therefore makes that partition
equal to the admissible lifecycle actions at the committed read state.
\end{proof}

This equality is revision-scoped: a lifecycle command re-evaluates its guards
before publication rather than treating an earlier Recall result as a durable
capability.

The complete branch work context also contains the structural neighborhood,
advisory dependency readiness, and branch-owned artifacts:
\begin{equation}
\Pi_b^{\mathrm{work}}(M_t,j)
=
\Bigl(
  \operatorname{RecallAPI}_t(b,j),\,
  N_t(b),\,
  \operatorname{Ready}_t(b),\,
  S_t(b),P_t(b),R_t(b),F_t(b),V_t(b)
\Bigr).
\label{eq:work-context-projection}
\end{equation}
Dependency readiness is deliberately advisory in the implemented system: it
guides branch choice but is not a hidden scheduler or a hard entry guard.  The
formal sufficiency claim concerns the declared project-level transaction guards
under $\Gamma$; readiness and the additional artifacts ground planning and
local implementation but do not guarantee code quality.

\begin{corollary}[Elimination of coordination-state aliasing]
\label{cor:no-coordination-aliasing}
For committed TreeWork states satisfying the assumptions of
Proposition~\ref{prop:treework-guard-complete}, two branch states with disjoint
admissible movement sets must produce different coordination projections
$\chi_t(b,j)$ and different decoded Recall eligibility.  Hence the
decisive-aliasing lower bound of
Theorem~\ref{thm:aliasing-lower-bound} does not arise from hidden
machine-owned coordination state once the matching committed Recall projection
is supplied.
\end{corollary}

\begin{corollary}[Cross-agent continuation agreement]
\label{cor:handoff-agreement}
If two agents with arbitrary private contexts read the same committed
Recall projection for the same invocation context, then they decode the same
admissible project-level continuation set for branch $b$.  They may still choose
different admissible actions or produce different code.
\end{corollary}

\subsection{Committed-State and Completion Soundness}

Let $\mathcal I(M)$ collect the structural and publication invariants of a
project memory, including one root, unique branch identifiers, a valid hierarchy,
an acyclic dependency graph, revision consistency, and event-tail consistency.

\begin{theorem}[Conditional committed-state soundness under marker-consistent publication]
\label{thm:committed-state-soundness}
Assume (i) $\mathcal I(M_0)$; (ii) for every successful transaction $\tau$ from
$M$ to $M'$, either $M'$ is explicitly validated against $\mathcal I$, or
$\tau$ is invariant-preserving,
\begin{equation}
\mathcal I(M)\land\operatorname{Guard}_\tau(M)
\Longrightarrow
\mathcal I(M');
\label{eq:invariant-preserving-transition}
\end{equation}
(iii) every interrupted or failed publication recovers either by restoring the
previous committed state and event tail or by finishing a complete intended
state that satisfies condition (ii); and (iv) marker-last publication defines
the logical commit point, while readers accept only marker-consistent tuples.
Then every accepted read is derived from one complete committed state $M_t$
satisfying $\mathcal I(M_t)$.
\end{theorem}

\begin{proof}
By induction on the committed transaction sequence, every committed state
satisfies $\mathcal I$: the base state does by assumption; rollback restores
the preceding committed state; finish-forward publishes the complete intended
state; an explicitly validated transaction publishes a state satisfying
$\mathcal I$; and an invariant-preserving transaction carries $\mathcal I$
forward under its guard.  Marker-last publication and marker-consistent reads
prevent an accepted reader from combining components belonging to different
logical commits.  Therefore every accepted read corresponds to one complete
committed state satisfying $\mathcal I$.
\end{proof}

This theorem is a conditional system invariant rather than an unconditional
claim about arbitrary files on disk.  The validation, journal, publication, and
coherent-read assumptions are connected to implementation mechanisms and tests
in Table~\ref{tab:runtime-correspondence}.

\begin{corollary}[Completion soundness]
\label{cor:completion-soundness}
If a committed transaction changes branch $b$ to
$\ell_{t+1}(b)=\textnormal{complete}$ only when
\begin{equation}
\operatorname{AcceptanceDone}_t(b)
\land
[v_t(b)=\textnormal{verified}]
\land
\operatorname{StateConsistent}_t
\land
\operatorname{WorkspaceSafe}_t(b),
\end{equation}
then every accepted completion satisfies all four declared completion
conditions at commit time.
\end{corollary}

\begin{proof}
The completion guard is evaluated before the committed lifecycle transition;
the premise requires every accepted transition to satisfy all four conjuncts.
\end{proof}
